\subsection{Identical finite metrics with different limiting velocity profiles}
\label{sec:different-velocity-same-metrics}
In the polarized class a common-node pair can also change
the limiting speed map itself while preserving both
induced metrics. In contrast to the rank-two example of
Section~\ref{sec:same-speed-different-modulus}, the
metric ambiguity here includes the leading local speed.

Fix regular times $a,b$ and let
$\mathcal K=\{k\ge1:u_k(b)=0\}$, using
\eqref{eq:invisible}. Consider polarized solutions
with prescribed constant $P(a),P(b)$ and $Q=Q_0$,
with $\lambda(a)$ constant fixed.
Then every compatible $P$ has the form
\begin{equation}
 P(t,\theta)=P(a)+\mu\log(t/a)+
 \sum_{k\in\mathcal K}u_k(t)
       (A_k\cos(k\theta)+B_k\sin(k\theta)),
 \quad \mu=\frac{P(b)-P(a)}{\log(b/a)}.
 \label{eq:polarized-fibre}
\end{equation}
This is the Fourier solution of the scalar equation
with the two prescribed constant displacements.
At both nodes the momentum constraint makes
$\lambda(a),\lambda(b)$ spatially constant.
The mean time quadrature gives
\begin{equation}
 \lambda(a)-\lambda(b)
  =\mu^2\log(b/a)+
     \frac12\sum_{k\in\mathcal K}(A_k^2+B_k^2)
         \int_a^b t\big(u_k'(t)^2+k^2u_k(t)^2\big)\,dt .
 \label{eq:fibre-energy}
\end{equation}
All the weights in the sum are positive.
This describes the full-metric fibre for
these constant observed orbit fields, with no
mode hidden in the zero-frequency part.
In particular, adding the terminal metric
fixes one weighted quadratic radius.

\begin{proposition}
\label{prop:different-velocity-same-metrics}
There exist polarized vacuum three-torus pairs with
identical induced spatial metrics and $Q_t$ at two
regular slices, but different angular profiles of
asymptotic speed. Both profiles can be bounded
strictly between zero and one everywhere.
The profiles can be attained on non-equivalent
histories under torus-equivariant isometries
preserving the observed boundary slices.
The construction is available for a dense set of
areal-time ratios $b/a>1$.
\end{proposition}

\begin{proof}
Choose two distinct modes in $\mathcal K$ with
$J_0(k a)\ne0$ for both. The modes $1,6$
from Proposition~\ref{prop:same-speed-different-modulus}
are one choice.
Use just one mode in each branch of
\eqref{eq:polarized-fibre}, matching the two
positive energies in \eqref{eq:fibre-energy}
by rescaling their nonzero amplitudes.
Here both $Q_t$ traces vanish since $Q$ is
constant. For $0<\mu<1$ and sufficiently
small amplitudes the two limiting profiles are
$\mu+\epsilon_j J_0(k_j a)\cos(k_j\theta)$
with values in $(0,1)$.
The distinct Fourier frequencies make them different.
For a geometric distinction already at the initial
boundary, its quotient metric is constant, so a
torus-equivariant isometry induces a rotation or
reflection of its circle. The mean curvature
depends there on $P_t^2$ with a common nonzero
factor and common additive term, since the
lapse and $\lambda(a)$ agree.
With nonzero $\mu$, the smallest nonzero harmonic
of this quantity is $k_j$ in branch $j$,
which no circle isometry can change.

For density, the ratio $\nu_5(\mathcal R)/
\nu_1(\mathcal R)$ of annular Dirichlet frequencies
is real analytic and nonconstant on $(1,\infty)$.
The limits at its endpoints, $5$ and
$j_{0,5}/j_{0,1}$, differ.
On each open interval of $\mathcal R$ its image
therefore contains an interval.
Choose a rational value $k_2/k_1$ in that image
and put $a=\nu_1(\mathcal R)/k_1$.
The two modes have common nodes.
The values $J_0(\nu_1(\mathcal R))$ and
$J_0(\nu_5(\mathcal R))$ vanish at most on
discrete sets of $\mathcal R$: each eigenvalue
is analytic and nonconstant.
Restrict the interval first to avoid those
discrete sets. This gives the required dense
collection of placements.
\end{proof}
